\section{Background}
Two commercial operators share a congested low-altitude coordination environment. Each knows its own confidentiality-cost type, but not the rival's. They choose Minimum, Partial, or Full voluntary commercial disclosure; Minimum still includes the model's assumed safety-information floor. Information sharing can improve routing coordination \cite{los2020value}, while competing operators and airspace institutions create strategic and privacy constraints \cite{liu2022efficient,maheshwari2025privacy}.

\section{Motivation}
The institutional problem is not “maximize disclosure.” Operators internalize confidentiality cost but jointly create congestion-dependent coordination benefits. A privately rational equilibrium can therefore miss total-welfare opportunities, while an access incentive can also induce excessive disclosure. The project asks where a candidate rule helps and where it fails, rather than treating one favorable calibration as a theorem.

\section{Application Scenario}
Operator A and B coordinate commercial flight intent in a capacity-constrained urban setting. Partial may encode a coarse next sector/ETA; Full is richer route/timing detail. Mandatory safety information, emergency missions, enforcement, inaccurate reports, and real traffic control lie outside the action space. Stakeholders include operators, a coordinator, customers, and affected publics. The model is a normalized computational benchmark, not an empirical aviation forecast.

\section{Research Question}
\textbf{Shared:} How do private confidentiality costs and public congestion shape voluntary disclosure, welfare, and the performance of a congestion-contingent access rule? \textbf{Economics} defines incentives and welfare; \textbf{computer science} exhaustively enumerates equilibria and falsifying parameter regions; \textbf{behavioral science} motivates future tests of comprehension, privacy salience, and fairness---no behavioral population result is claimed.

\section{Method}
Nature independently draws $c_i\in\{1,1.6\}$ with probability $1/2$; public congestion is $d\in\{0.3,0.6,1\}$. Simultaneous actions are $a_i\in\{M,P,F\}$ with $x=(0,1,2)$ and $k=(0,1,2.5)$. Utility is
\[u_i=d(7.5)[1-e^{-0.5(x_i+x_j)}]-c_i k(a_i).\]
A complete pure strategy maps both own types to actions. The solver checks all $9\times9=81$ strategy pairs for interim best responses; symmetry is not imposed. Welfare is $W=u_A+u_B$. The full-information first best observes realized types and is a reference, not necessarily feasible. CCAR scales only a Minimum discloser's access to enhanced, non-safety coordination benefit by $\alpha_M=0.7$ at Medium/High congestion. It does not withhold mandatory safety information or establish truthful reporting. The downstream allocation appendix separately compares FCFS with a one-slot second-price auction; it does not turn the disclosure game into an auction.

\textbf{Artifacts:} \PublicRepo; \ColabOne; \ColabTwo; \SpaceRepo. Validated source commit: \TestedCommit.

\balance
\section{Results}
All benchmark states have one pure BNE. Low: $(M,M)$, $W=0.000$, first best $0.639$, gap $0.639$. Medium: $(P,M)$ by type, $W=2.193$, first best $3.139$, gap $0.946$. High: $(P,P)$, $W=6.882$, first best $7.385$, gap $0.503$. Thus the largest benchmark gap is Medium. There, CCAR changes the BNE to $(P,P)$ and underlying welfare to $3.089$, leaving gap $0.050$. These are code-recomputed normalized expectations. Across 760 deterministic grid cells, diagnostics overlap: 401 no equilibrium-set effect, 334 robust welfare improvements, 18 excessive-disclosure flags, 3 larger-gap cases, 66 multiple-CCAR-BNE cases, and 0 without a pure CCAR BNE. The cells are neither subjects nor probabilities.

\section{Intellectual Merits}
Los et al. quantify operational sharing value in collaborative vehicle routing \cite{los2020value}; Liu and Chow study efficient/stable transit-operator data sharing \cite{liu2022efficient}; Maheshwari et al. design privacy-preserving AAM coordination \cite{maheshwari2025privacy}. FP10's bounded difference is the exact link among private binary confidentiality types, three disclosure levels, public congestion, complete Bayesian strategies, a realized-type first-best benchmark, and CCAR failure regions. This is a reproducible application/model combination and comparative numerical result---not a general theorem, proof of optimality, or empirical novelty claim. Appendix G expands Figure~\ref{fig:teaser}.

\section{Practical Impacts}
Potential relevance is limited to UN SDG targets 9.1 (sustainable/resilient infrastructure) and 11.2 (safe, accessible, sustainable transport). The demonstrated evidence is only stylized welfare/disclosure computation. The SDG 11.2 link is indirect and prospective: this model does not evaluate passenger access, public-transport accessibility, or field safety. Proposed indicators are coordination delay, confidential-data exposure, access disparity, and override/safety incidents relative to an unrestricted-access baseline. Risks include high-cost-operator disadvantage, gaming, and over-disclosure; safeguards are a safety-information floor, auditable eligibility, appeal, privacy review, and suspension on adverse diagnostics. Academia requires reproducibility; industry additionally requires reliability/security; public institutions require legitimacy, equitable access, contestability, and correction. No emissions or real-world safety benefit is claimed.

\subsection*{Future Research Roadmap}
Economics asks whether reporting and participation constraints remain feasible with implementation cost; a mechanism-design or transport-economics collaborator would test incentive compatibility and individual rationality across documented parameter regions, culminating in a reproducible mechanism extension. Computer science asks whether CCAR remains safe and useful in dynamic heterogeneous traffic; a multi-agent, transport-simulation, or privacy-computing collaborator would report runtime, throughput, safety-constraint, and privacy-leakage metrics, culminating in a reproducible simulation benchmark. Behavioral science asks whether people understand the incentives and fairness implications; a behavioral economist, HCI researcher, or experimental scientist would preregister a protocol, comprehension check, and choice measure before an ethically reviewed pilot. Appendix~G, Figure~4 labels each question, method/data source, potential collaborator, progress criterion, and milestone. Han Zhang's suggested independent-coordinator role remains a proposed extension, not a completed model or collaboration.

\subsection*{Evolving 2056 Contribution Statement}
By 2056, the aspiration is accountable incentive-aware coordination in shared autonomous systems. The current contribution is only the verified stylized comparison above. Official records identify the 2025 prize to Joel Mokyr, Philippe Aghion, and Peter Howitt for explaining innovation-driven economic growth \cite{nobel2025growth}, and the 2026 prize to Robert E. Hall and Michael Woodford for microeconomic foundations of macroeconomic stabilisation \cite{kva2026stabilisation}. These ideas motivate long-horizon mechanism-building grounded in agent behavior; neither prize validates FP10. Progress requires general theory, independent replication, behavioral evidence, and safe field evaluation.

\subsection*{Open Science and SDG Statements}
Shortest local reproduction runs the repository's central evidence builder; model data, grids, figures, tests, and manifests are public alongside it. Local notebook execution is recorded; hosted Colab and final deployment parity remain separately labeled. SDG connections are prospective until stakeholder-defined field indicators and a governance process exist.
